\subsection{Informes durables de expedientes}\label{sec:impl-informes}

La implementación local compone la solicitud autenticada, la persistencia
transaccional, el cifrado de capturas y artefactos, el renderizado acotado y la
consulta desde Qadra. \texttt{CaseReportWorkflow} delimita el puerto de aplicación;
HTTP expone solicitud, listado propio, detalle, selector de litigantes, acuse y
descarga bajo \texttt{/api/v1/case-reports}. El adaptador PostgreSQL comprueba la
identidad completa, la generación de autorización y las pertenencias dentro de
la transacción auditada. Los cambios posteriores de permisos no quedan
neutralizados por conservar una sesión o por devolver a la cuenta su rol
anterior. El consumo de recuperación que sólo incrementa la revisión de cuenta
no invalida por sí mismo el informe.

El intervalo filtra la creación original y el estado se observa al capturar,
según la \cref{sec:diseno-informes}. El filtro opcional por litigante limita los
expedientes; no elimina del resumen a otros litigantes asignados a esos mismos
casos. La captura conserva filtros, instante de observación, identificadores,
revisiones administrativas, asignaciones y carga derivada. Repetir una operación
con sus mismos filtros devuelve el trabajo original; cambiar los filtros con
el mismo identificador produce conflicto. La interfaz conserva esa identidad
ante una respuesta incierta y permite reintentar explícitamente.

\paragraph{Persistencia y recuperación.}
La migración \rutarepo{migrations/0026_case_reports.sql} añade solicitudes,
capturas, artefactos y avisos. Las solicitudes mantienen la identidad original
junto con el estado de ejecución. Capturas y artefactos son inmutables: se
protegen mediante los puertos de cifrado autenticado y envoltura de claves,
ligando identificador de informe, tipo de contenido y digest. Los dos formatos
proceden de los mismos bytes canónicos, incluso después de reintentar un
renderizado. No se crea una identidad documental ni un sello jurídico para la
exportación. La KEK y la copia consistente de PostgreSQL son necesarias para
restaurar solicitudes, autorizaciones y resultados.

\texttt{serve} incorpora un consumidor serial supervisado con parada coordinada.
La cola persistente concede cada trabajo por dos minutos mediante identificador,
intento, token y generación; las renovaciones amplían la concesión en dos
minutos. Cada escritura comprueba la concesión completa y su vigencia. Un
proceso anterior no puede publicar sobre un intento posterior. Los errores
transitorios elegibles admiten hasta cinco intentos, con una espera de treinta
segundos multiplicada por el número de intento. El procesamiento puede recuperar
un trabajo interrumpido sin recapturar un resultado ya persistido.

El renderizado se ejecuta fuera de la transacción de auditoría. Antes de publicar
se revalida el acceso a toda la captura y una sola confirmación guarda ambos
artefactos, el estado disponible y el aviso. También la descarga comprueba el
conjunto completo. Perder una pertenencia impide obtener el archivo, sin
recortarlo ni fabricar una nueva captura. El acuse registra únicamente la
primera lectura del aviso; consultar el detalle o descargar no lo sustituye.

\paragraph{Presupuestos del renderizador.}
El adaptador Linux invoca una entrada privada del mismo ejecutable confiable por
cada formato. Recibe la captura en memoria sellada, con entorno vacío, sin los
descriptores heredados de los servicios y con un protocolo de salida ligado al
informe y al digest. El proceso padre controla el grupo hijo y su terminación.
Los límites por formato son quince segundos de CPU, 512 MiB de espacio virtual
y veinte segundos de tiempo de pared. El crecimiento de archivos regulares y
los volcados de memoria tienen límite cero; la salida se acota al artefacto de
16 MiB y su cabecera. Estos controles limitan recursos y separan la ejecución,
pero no crean un aislamiento adicional de permisos, archivos o red.

La captura admite hasta mil expedientes, diez mil pares de asignación y mil
litigantes en el resumen, con un máximo canónico de 8 MiB. El PDF limita además
la composición a 512 páginas y un millón de glifos, con fuentes incorporadas y
rechazo explícito de caracteres no soportados. El CSV utiliza registros RFC
4180 y protege textos no confiables frente a fórmulas de hoja de cálculo. Los
presupuestos se aplican conjuntamente: el máximo de filas no garantiza que
el resultado quepa en los límites de páginas, memoria y tiempo, incluso con
campos cortos. Un exceso falla sin omitir datos para aparentar un informe completo.
Cada instancia de fuente reutiliza su plan tipográfico para medir y dibujar
texto latino de izquierda a derecha con las mismas opciones. Las líneas se
componen por separado para respetar sus extremos y sus ligaduras;
no se rebajan los presupuestos para aceptar un artefacto.

\paragraph{Presentación y alcance de la entrega.}
Qadra permite filtrar, solicitar, actualizar el estado, consultar trabajos,
acusar avisos y descargar un formato. Distingue espera, captura, renderizado,
reintento, disponibilidad y fallo; no estima porcentajes ni fechas de finalización.
El selector paginado conserva los colegas compartidos sólo en expedientes
cerrados. Antes de ofrecer una descarga, el cliente contrasta identificador,
captura, tamaño y digest con el detalle autorizado.

El contrato \rutarepo{docs/case-reports-api.md} y la operación de respaldo en
\rutarepo{docs/database-operations.md} describen estas fronteras. La decisión
\rutarepo{docs/adr/0060-durable-authorized-case-reports.md} documenta las decisiones de autorización,
persistencia y publicación conjunta verificadas en la aceptación integrada. La composición local
no acredita por sí sola un despliegue, una campaña global aprobada, efectividad
jurídica, firma personal o constancia emitida por un PSC. El aviso por correo y
la estimación de terminación de la \cref{tab:cu-16}, junto con otros tipos de
informe y las métricas de desempeño de la \cref{tab:rf-19}, conservan su alcance
pendiente. La exportación de estado y carga no completa esos criterios.
